\scanpage{25}
\mspage{24}
\mspage{IV-950}

$f_0^{(\lambda)}:S_0^{(\lambda)}\to S_0'{}^{(\lambda)}$ tels que $S'=S'{}^{(\lambda)}\times_{S^{(\lambda)}}S$, $\unsure{g|f}=\unsure{g|f}^{(\lambda)}\times1$ et $f_0=f_0^{(\lambda)}\times1$. \deleted{\illegible{1}} Soient $s$ le point fermé de $S$, $x=f_0(s)\unsure{_s}$, $x_\lambda$ la projection de $x$ dans $S'{}^{(\lambda)}$; comme $x_\lambda$ est au-dessus du point fermé de $S^{(\lambda)}$, l'anneau local \inserted{$\mathcal O_\lambda$} de $S'{}^{(\lambda)}$ au point $x_\lambda$ est une $A_\lambda$-algèbre essentiellement étale; en outre, comme $f_0^{(\lambda)}(s_\lambda)=$ $=x_\lambda$, on a $\kappa(x_\lambda)=\kappa(s_\lambda)$, autrement dit, $\mathcal O_\lambda$ est une $A_\lambda$-algèbre locale strictement essentiellement étale, et par suite (18.6.1) est $A_\lambda$-isomorphe à une $A$-algèbre $A_\mu$ avec $\mu\geq\lambda$. Il y a \deleted{\illegible{2}}\inserted{donc} un $A_\lambda$-homomorphisme \deleted{$\unsure{A_\mu\to A}$}\inserted{$\mathcal O_\lambda\to\unsure{A_\lambda|A}$}, c'est-à-dire un $S^{(\lambda)}$-morphisme $h:S\to S'{}^{(\lambda)}$ tel que $h(s)=x_\lambda$, et par suite il existe bien une $S$-section $f$ de $S'$ telle que $f(s)=x$, ce qui achève de prouver (i);

Pour prouver (ii), il suffit de noter que l'on a $\operatorname{Hom.loc}({}^h A,B)\deleted{\unsure{=}}$ $\deleted{=}\inserted{\simeq}\varprojlim\operatorname{Hom.loc}(A_\lambda,B)$ et que par (18.6.2) les homomorphismes canoniques $\operatorname{Hom.loc}(A_\lambda,B)\leftarrow\operatorname{Hom.loc}(A,B)$ sont bijectifs.

\deleted{\unsure{Enfin}} Pour démontrer \deleted{\unsure{le}} (iv), \deleted{\unsure{remarquons}} notons que pour tout $\lambda$ et tout entier $n>0$, $\underline m_\lambda^n=\underline m^n A_\lambda$, et $(\underline m.{}^h A)^n=\underline m^n.{}^h A=\varinjlim\underline m_\lambda^n$. On en conclut, par l'exactitude du foncteur $\varinjlim$ dans la catégorie des $A$-modules, que ${}^h A/(\underline m.{}^h A)^n=$ $=\varinjlim(A_\lambda/\underline m_\lambda^n)$, et par suite, il suffit de montrer que, pour tout entier $n$, et tout indice $\lambda$, l'homomorphisme $A/\underline m^n\to A_\lambda/\underline m_\lambda^n$ est bijectif. Or, \deleted{\unsure{pour}} cela est vrai par hypothèse pour $n=1$; d'autre part, puisque $A_\lambda$ est un $A$-module plat, on a
\[
\underline m_\lambda^n/\underline m_\lambda^{n+1}=(\underline m^n/\underline m^{n+1})\otimes_A A_\lambda=(\underline m^n/\underline m^{n+1})\otimes_{A/\underline m}(A_\lambda/\underline m A_\lambda)
\]
et comme $A/\underline m\to A_\lambda/\underline m A_\lambda=A_\lambda/\underline m_\lambda$ est bijectif, l'homomorphisme $\underline m^n/\underline m^{n+1}\to\underline m_\lambda^n/\underline m_\lambda^{n+1}$ est lui aussi bijectif; la conclusion résulte donc de Bourbaki, \uline{Alg.comm.}, chap.III,§ 2,n°8,\inserted{cor.3 du th.1}.

Enfin, si $A$ est hensélien \deleted{[et si $A'=B_{\underline n}$ est une $A$-algèbre locale strictement essentiellement étale, où $B$ est une $A$-algèbre étale, il résulte de \deleted{\unsure{la}} (18.6.2) qu'il existe une \deleted{\unsure{unique}} \uline{$\operatorname{Spec}(A)$-section de} $\operatorname{Spec}(B)$ prenant la valeur $\underline n$ au point fermé de $\operatorname{Spec}(A)$]}, on conclut \deleted{alors de}\inserted{de la remarque qui suit} (\deleted{17.4.1}\inserted{18.6.2}) que \deleted{\unsure{nécessairement}} \deleted{l'}\inserted{les}homomorphisme\inserted{s} \deleted{\unsure{canonique}} $A\to A\deleted{'}\inserted{_\lambda}$ \deleted{est}\inserted{sont} bijectif, ce qui prouve (vi) par définition de ${}^h A$.

(18.6.7) Soient maintenant $A$ un anneau \uline{semi-local}, $\underline m_{\unsure{1|i}}$ $(1\leq\unsure{1|i}\leq r)$ ses idéaux maximaux. On appelle \uline{hensélisé} de $A$ et l'on note encore ${}^h A$ l'anneau \uline{produit} $\prod_i{}^h(A_{\underline m_{\unsure{1|i}}})$ des hensélisés des anneaux locaux $A_{\underline m_{\unsure{1|i}}}$. C'est \inserted{un $A$-module fidèlement plat et} une $A$-algèbre semi-locale, dont les idéaux maximaux sont les $\underline m_{\unsure{1|i}}.{}^h A$ en vertu de (18.6.6,(iii)); en outre, si $\underline r=\bigcap_i\underline m_{\unsure{1|i}}$ est le radical de $A$, il résulte de ce qui précède que $\underline r.{}^h A$ est le radical de ${}^h A$ et que l'application $A/\underline r\to{}^h A/\underline r.{}^h A$ est bijective. Comme le séparé complété $\widehat A$ de $A$ pour la topologie $\underline r$-préadique est le produit des séparés complétés $\widehat{A_{\underline m_{\unsure{1|i}}}}$ (Bourbaki,\uline{Alg.comm.},chap.III,§ 2,n°13\inserted{prop.18}) l'homomorphisme canonique $\widehat A\to({}^h A)^\wedge$ est bijectif par (18.6.6,(iv)).
